\section*{ID9831: Euler–Lagrange Equations and Gradient Conditions: (∂ J)/(∂ θ\_j)(bmθ\textasciicircum{}*)}
\paragraph{Metadata.}
PiID: 6513410559240 | RTEID: RTE:P38175.9562:T0.0704 | UUID: d8264a9b-bcdc-522e-a298-ad6b31bcbaea

\paragraph{Canonical Statement.}
Formally, the stationary points of \$\textbackslash{}mathcal\{J\}\$ satisfy \textbackslash{}[ \textbackslash{}frac\{\textbackslash{}partial \textbackslash{}mathcal\{J\}\}\{\textbackslash{}partial \textbackslash{}theta\_j\}(\textbackslash{}bm\{\textbackslash{}theta\}\textasciicircum{}*) = 0 \textbackslash{}quad \textbackslash{}text\{for all components \} \textbackslash{}theta\_j \textbackslash{}text\{ of \} \textbackslash{}bm\{\textbackslash{}theta\}. \textbackslash{}]

\paragraph{Canonical Equation.}
\[
\frac{\partial \mathcal{J}}{\partial \theta_j}(\bm{\theta}^*) = 0 \quad \text{for all components } \theta_j \text{ of } \bm{\theta}.
\]

\paragraph{Dictionary.}
Euler–Lagrange Equations and Gradient Conditions: (∂ J)/(∂ θ\_j)(bmθ\textasciicircum{}*) — (∂ J)/(∂ θ\_j)(bmθ\textasciicircum{}*) = 0 for all components θ\_j of bmθ.

\paragraph{Encyclopedia.}
Formally, the stationary points of satisfy It arises from a variational or Legendre procedure, defining conjugate quantities or stationarity conditions.
